% --- Archon physics formalization source begin ---
% archon:physics
% archon:covers PhyXMiniProblems/problem_phyx_mini_0930.lean
% archon:source-report reports/phyx_mini/problem_phyx_mini_0930.source.json

\chapter{Physics problem phyx\_mini\_0930}
\label{ch:PhyXMiniProblems_problem_phyx_mini_0930}

\paragraph{Problem source.}
\#\# Physical scenario

Figure shows a \textbackslash{}( 4.0 \textbackslash{})-cm-diameter loop with resistance \textbackslash{}( 0.10 \textbackslash{}, \textbackslash{}Omega \textbackslash{}) around a \textbackslash{}( 2.0 \textbackslash{})-cm-diameter solenoid. The solenoid is \textbackslash{}( 10 \textbackslash{}, \textbackslash{}text\{cm\} \textbackslash{}) long, has \textbackslash{}( 100 \textbackslash{}) turns, and carries the current shown in the graph. A positive current is cw when seen from the left.

\#\# Question

Find the current in the loop at \textbackslash{}( t = 2.5 \textbackslash{}, \textbackslash{}text\{s\} \textbackslash{}).

\#\# Answer choices

- A: A: \textbackslash{}( 1.6 \textbackslash{} \textbackslash{}mu \textbackslash{}text\{A\} \textbackslash{})

- B: B: \textbackslash{}( 0 \textbackslash{} \textbackslash{}mu \textbackslash{}text\{A\} \textbackslash{})

- C: C: \textbackslash{}( 3.2 \textbackslash{} \textbackslash{}mu \textbackslash{}text\{A\} \textbackslash{})

- D: D: \textbackslash{}( 0.8 \textbackslash{} \textbackslash{}mu \textbackslash{}text\{A\} \textbackslash{})

\#\# Figure caption (auxiliary; use the image as primary evidence)

The image consists of two parts:

1. **Left Diagram**: 
   - **Object**: A side view of a solenoid is shown.
   - **Description**: The solenoid is depicted with multiple loops wrapped around a cylindrical core, shown in cross-section. The loops are represented as purple circles on the top and bottom sides.
   - **Arrow**: There is an arrow labeled \textbackslash{}( I\_\{\textbackslash{}text\{sol\}\} \textbackslash{}) indicating the direction of current flow in the solenoid. 
   - **Dimensions**: The solenoid's height is labeled as 4.0 cm and its diameter as 2.0 cm.

2. **Right Graph**: 
   - **Axes**: The graph has time \textbackslash{}( t \textbackslash{}) in seconds on the x-axis, ranging from 0 to 3 seconds, and the solenoid current \textbackslash{}( I\_\{\textbackslash{}text\{sol\}\} \textbackslash{}) in amperes on the y-axis, ranging from -20 to 20 amperes.
   - **Plot**: A purple line graph shows the current behavior. 
     - From 0 to 1 second, the current is constant at 20 A.
     - From 1 to 2 seconds, the current decreases linearly to -20 A.
     - From 2 to 3 seconds, the current remains constant at -20 A.

This figure illustrates the configuration and current change over time for the solenoid.

\#\# Dataset metadata

Electromagnetic Induction; Multi-Formula Reasoning, Physical Model Grounding Reasoning

\paragraph{Recorded answer/context.}
B: B: \textbackslash{}( 0 \textbackslash{} \textbackslash{}mu \textbackslash{}text\{A\} \textbackslash{})

\paragraph{Figure/image path.}
/root/proposal\_for\_physic/science-mango/phyx\_mini\_run/phyx\_data/test\_image/930.png

\paragraph{Formalization target.}
create a compiling Lean file with sorry bodies at `PhyXMiniProblems/problem\_phyx\_mini\_0930.lean`. The Lean declarations must preserve the physical quantities, dimensions or dimensional roles, figure labels, governing-law hypotheses, and final relation expressed by this problem.
Use Mathlib/Physlib names found through LeanExplore where available. If a physics API is missing, introduce faithful local abstractions rather than scalar placeholder aliases.

\begin{theorem}[Physics formalization target]
\label{thm:physics:phyx_mini_0930:target}
\lean{PhyXMiniProblems.ProblemPhyXMini0930.problem_phyx_mini_0930}
\uses{def:physics:phyx-mini-0930:phyxminiproblems-problemphyxmini0930-solenoidloopsetup, def:physics:phyx-mini-0930:phyxminiproblems-problemphyxmini0930-matchesstatedsolenoidloopscenario, def:physics:phyx-mini-0930:phyxminiproblems-problemphyxmini0930-matchessuppliedsolenoidloopfigure, def:physics:phyx-mini-0930:phyxminiproblems-problemphyxmini0930-hasphysicalsolenoidloopparameters, def:physics:phyx-mini-0930:phyxminiproblems-problemphyxmini0930-satisfiescircularsolenoidgeometry, def:physics:phyx-mini-0930:phyxminiproblems-problemphyxmini0930-satisfiessolenoidcurrentratedefinition, def:physics:phyx-mini-0930:phyxminiproblems-problemphyxmini0930-satisfiesidealsolenoidfieldlaw, def:physics:phyx-mini-0930:phyxminiproblems-problemphyxmini0930-satisfiesuniformsolenoidfluxlinkagelaw, def:physics:phyx-mini-0930:phyxminiproblems-problemphyxmini0930-satisfiesfaradayinductionlaw, def:physics:phyx-mini-0930:phyxminiproblems-problemphyxmini0930-satisfiesresistiveloopohmslaw, def:physics:phyx-mini-0930:phyxminiproblems-problemphyxmini0930-recordeddatasetanswer, def:physics:phyx-mini-0930:phyxminiproblems-problemphyxmini0930-observedloopcurrentinmicroamperes, def:physics:phyx-mini-0930:phyxminiproblems-problemphyxmini0930-isuniquematchinganswer}
The assigned autoformalize agent should translate this physics problem into Lean declarations in the covered file, with theorem and lemma proof bodies written as `by sorry`.
\end{theorem}
\begin{proof}
This is an autoformalization task, not a proof task. Produce statements that can later be proved by the physics prover without weakening the physical model.
\end{proof}

% --- Archon declaration topology begin ---
\section{Declaration topology}
\begin{definition}[electric Current Dimension]
  \label{def:physics:phyx-mini-0930:phyxminiproblems-problemphyxmini0930-electriccurrentdimension}
  \lean{PhyXMiniProblems.ProblemPhyXMini0930.electricCurrentDimension}
  Electric current has dimension charge per time
\end{definition}
\begin{proof}
  This is the declared model definition of electric Current Dimension.
\end{proof}
\begin{definition}[electric Current Rate Dimension]
  \label{def:physics:phyx-mini-0930:phyxminiproblems-problemphyxmini0930-electriccurrentratedimension}
  \lean{PhyXMiniProblems.ProblemPhyXMini0930.electricCurrentRateDimension}
  \uses{def:physics:phyx-mini-0930:phyxminiproblems-problemphyxmini0930-electriccurrentdimension}
  The time derivative of current has dimension charge per time squared
\end{definition}
\begin{proof}
  This is the declared model definition of electric Current Rate Dimension.
\end{proof}
\begin{definition}[magnetic Permeability Dimension]
  \label{def:physics:phyx-mini-0930:phyxminiproblems-problemphyxmini0930-magneticpermeabilitydimension}
  \lean{PhyXMiniProblems.ProblemPhyXMini0930.magneticPermeabilityDimension}
  Magnetic permeability has dimension M L C⁻²
\end{definition}
\begin{proof}
  This is the declared model definition of magnetic Permeability Dimension.
\end{proof}
\begin{definition}[magnetic Flux Density Dimension]
  \label{def:physics:phyx-mini-0930:phyxminiproblems-problemphyxmini0930-magneticfluxdensitydimension}
  \lean{PhyXMiniProblems.ProblemPhyXMini0930.magneticFluxDensityDimension}
  Magnetic flux density has dimension M T⁻¹ C⁻¹
\end{definition}
\begin{proof}
  This is the declared model definition of magnetic Flux Density Dimension.
\end{proof}
\begin{definition}[magnetic Flux Density Rate Dimension]
  \label{def:physics:phyx-mini-0930:phyxminiproblems-problemphyxmini0930-magneticfluxdensityratedimension}
  \lean{PhyXMiniProblems.ProblemPhyXMini0930.magneticFluxDensityRateDimension}
  \uses{def:physics:phyx-mini-0930:phyxminiproblems-problemphyxmini0930-magneticfluxdensitydimension}
  A magnetic-flux-density change rate has one additional inverse time
\end{definition}
\begin{proof}
  This is the declared model definition of magnetic Flux Density Rate Dimension.
\end{proof}
\begin{definition}[area Dimension]
  \label{def:physics:phyx-mini-0930:phyxminiproblems-problemphyxmini0930-areadimension}
  \lean{PhyXMiniProblems.ProblemPhyXMini0930.areaDimension}
  Area has dimension length squared
\end{definition}
\begin{proof}
  This is the declared model definition of area Dimension.
\end{proof}
\begin{definition}[magnetic Flux Dimension]
  \label{def:physics:phyx-mini-0930:phyxminiproblems-problemphyxmini0930-magneticfluxdimension}
  \lean{PhyXMiniProblems.ProblemPhyXMini0930.magneticFluxDimension}
  \uses{def:physics:phyx-mini-0930:phyxminiproblems-problemphyxmini0930-magneticfluxdensitydimension, def:physics:phyx-mini-0930:phyxminiproblems-problemphyxmini0930-areadimension}
  Magnetic flux is magnetic flux density times area
\end{definition}
\begin{proof}
  This is the declared model definition of magnetic Flux Dimension.
\end{proof}
\begin{definition}[electromotive Force Dimension]
  \label{def:physics:phyx-mini-0930:phyxminiproblems-problemphyxmini0930-electromotiveforcedimension}
  \lean{PhyXMiniProblems.ProblemPhyXMini0930.electromotiveForceDimension}
  \uses{def:physics:phyx-mini-0930:phyxminiproblems-problemphyxmini0930-magneticfluxdimension}
  Electromotive force has the same dimension as magnetic-flux change rate
\end{definition}
\begin{proof}
  This is the declared model definition of electromotive Force Dimension.
\end{proof}
\begin{definition}[electrical Resistance Dimension]
  \label{def:physics:phyx-mini-0930:phyxminiproblems-problemphyxmini0930-electricalresistancedimension}
  \lean{PhyXMiniProblems.ProblemPhyXMini0930.electricalResistanceDimension}
  \uses{def:physics:phyx-mini-0930:phyxminiproblems-problemphyxmini0930-electriccurrentdimension, def:physics:phyx-mini-0930:phyxminiproblems-problemphyxmini0930-electromotiveforcedimension}
  Electrical resistance has the dimension emf divided by current
\end{definition}
\begin{proof}
  This is the declared model definition of electrical Resistance Dimension.
\end{proof}
\begin{definition}[Length Magnitude]
  \label{def:physics:phyx-mini-0930:phyxminiproblems-problemphyxmini0930-lengthmagnitude}
  \lean{PhyXMiniProblems.ProblemPhyXMini0930.LengthMagnitude}
  A nonnegative, unit-independent physical length
\end{definition}
\begin{proof}
  This is the declared model definition of Length Magnitude.
\end{proof}
\begin{definition}[Time Magnitude]
  \label{def:physics:phyx-mini-0930:phyxminiproblems-problemphyxmini0930-timemagnitude}
  \lean{PhyXMiniProblems.ProblemPhyXMini0930.TimeMagnitude}
  A nonnegative, unit-independent physical time
\end{definition}
\begin{proof}
  This is the declared model definition of Time Magnitude.
\end{proof}
\begin{definition}[Area Magnitude]
  \label{def:physics:phyx-mini-0930:phyxminiproblems-problemphyxmini0930-areamagnitude}
  \lean{PhyXMiniProblems.ProblemPhyXMini0930.AreaMagnitude}
  \uses{def:physics:phyx-mini-0930:phyxminiproblems-problemphyxmini0930-areadimension}
  A nonnegative, unit-independent physical area
\end{definition}
\begin{proof}
  This is the declared model definition of Area Magnitude.
\end{proof}
\begin{definition}[Signed Electric Current]
  \label{def:physics:phyx-mini-0930:phyxminiproblems-problemphyxmini0930-signedelectriccurrent}
  \lean{PhyXMiniProblems.ProblemPhyXMini0930.SignedElectricCurrent}
  \uses{def:physics:phyx-mini-0930:phyxminiproblems-problemphyxmini0930-electriccurrentdimension}
  A signed electric current relative to a stated circuit orientation
\end{definition}
\begin{proof}
  This is the declared model definition of Signed Electric Current.
\end{proof}
\begin{definition}[Signed Electric Current Rate]
  \label{def:physics:phyx-mini-0930:phyxminiproblems-problemphyxmini0930-signedelectriccurrentrate}
  \lean{PhyXMiniProblems.ProblemPhyXMini0930.SignedElectricCurrentRate}
  \uses{def:physics:phyx-mini-0930:phyxminiproblems-problemphyxmini0930-electriccurrentratedimension}
  A signed time derivative of electric current
\end{definition}
\begin{proof}
  This is the declared model definition of Signed Electric Current Rate.
\end{proof}
\begin{definition}[Magnetic Permeability Magnitude]
  \label{def:physics:phyx-mini-0930:phyxminiproblems-problemphyxmini0930-magneticpermeabilitymagnitude}
  \lean{PhyXMiniProblems.ProblemPhyXMini0930.MagneticPermeabilityMagnitude}
  \uses{def:physics:phyx-mini-0930:phyxminiproblems-problemphyxmini0930-magneticpermeabilitydimension}
  A nonnegative magnetic permeability
\end{definition}
\begin{proof}
  This is the declared model definition of Magnetic Permeability Magnitude.
\end{proof}
\begin{definition}[Signed Magnetic Flux Density]
  \label{def:physics:phyx-mini-0930:phyxminiproblems-problemphyxmini0930-signedmagneticfluxdensity}
  \lean{PhyXMiniProblems.ProblemPhyXMini0930.SignedMagneticFluxDensity}
  \uses{def:physics:phyx-mini-0930:phyxminiproblems-problemphyxmini0930-magneticfluxdensitydimension}
  A signed axial component of magnetic flux density
\end{definition}
\begin{proof}
  This is the declared model definition of Signed Magnetic Flux Density.
\end{proof}
\begin{definition}[Signed Magnetic Flux Density Rate]
  \label{def:physics:phyx-mini-0930:phyxminiproblems-problemphyxmini0930-signedmagneticfluxdensityrate}
  \lean{PhyXMiniProblems.ProblemPhyXMini0930.SignedMagneticFluxDensityRate}
  \uses{def:physics:phyx-mini-0930:phyxminiproblems-problemphyxmini0930-magneticfluxdensityratedimension}
  A signed time derivative of an axial magnetic-flux-density component
\end{definition}
\begin{proof}
  This is the declared model definition of Signed Magnetic Flux Density Rate.
\end{proof}
\begin{definition}[Signed Magnetic Flux]
  \label{def:physics:phyx-mini-0930:phyxminiproblems-problemphyxmini0930-signedmagneticflux}
  \lean{PhyXMiniProblems.ProblemPhyXMini0930.SignedMagneticFlux}
  \uses{def:physics:phyx-mini-0930:phyxminiproblems-problemphyxmini0930-magneticfluxdimension}
  Signed magnetic flux relative to the selected surface normal
\end{definition}
\begin{proof}
  This is the declared model definition of Signed Magnetic Flux.
\end{proof}
\begin{definition}[Signed Magnetic Flux Rate]
  \label{def:physics:phyx-mini-0930:phyxminiproblems-problemphyxmini0930-signedmagneticfluxrate}
  \lean{PhyXMiniProblems.ProblemPhyXMini0930.SignedMagneticFluxRate}
  \uses{def:physics:phyx-mini-0930:phyxminiproblems-problemphyxmini0930-electromotiveforcedimension}
  Signed magnetic-flux change rate
\end{definition}
\begin{proof}
  This is the declared model definition of Signed Magnetic Flux Rate.
\end{proof}
\begin{definition}[Signed Electromotive Force]
  \label{def:physics:phyx-mini-0930:phyxminiproblems-problemphyxmini0930-signedelectromotiveforce}
  \lean{PhyXMiniProblems.ProblemPhyXMini0930.SignedElectromotiveForce}
  \uses{def:physics:phyx-mini-0930:phyxminiproblems-problemphyxmini0930-electromotiveforcedimension}
  Signed electromotive force relative to the positive loop orientation
\end{definition}
\begin{proof}
  This is the declared model definition of Signed Electromotive Force.
\end{proof}
\begin{definition}[Electrical Resistance Magnitude]
  \label{def:physics:phyx-mini-0930:phyxminiproblems-problemphyxmini0930-electricalresistancemagnitude}
  \lean{PhyXMiniProblems.ProblemPhyXMini0930.ElectricalResistanceMagnitude}
  \uses{def:physics:phyx-mini-0930:phyxminiproblems-problemphyxmini0930-electricalresistancedimension}
  A nonnegative electrical resistance
\end{definition}
\begin{proof}
  This is the declared model definition of Electrical Resistance Magnitude.
\end{proof}
\begin{definition}[signed SIReadout]
  \label{def:physics:phyx-mini-0930:phyxminiproblems-problemphyxmini0930-signedsireadout}
  \lean{PhyXMiniProblems.ProblemPhyXMini0930.signedSIReadout}
  Coherent-SI readout of a signed dimensionful quantity
\end{definition}
\begin{proof}
  This is the declared model definition of signed SIReadout.
\end{proof}
\begin{definition}[nonnegative SIReadout]
  \label{def:physics:phyx-mini-0930:phyxminiproblems-problemphyxmini0930-nonnegativesireadout}
  \lean{PhyXMiniProblems.ProblemPhyXMini0930.nonnegativeSIReadout}
  Coherent-SI readout of a nonnegative dimensionful quantity
\end{definition}
\begin{proof}
  This is the declared model definition of nonnegative SIReadout.
\end{proof}
\begin{definition}[length In Meters]
  \label{def:physics:phyx-mini-0930:phyxminiproblems-problemphyxmini0930-lengthinmeters}
  \lean{PhyXMiniProblems.ProblemPhyXMini0930.lengthInMeters}
  \uses{def:physics:phyx-mini-0930:phyxminiproblems-problemphyxmini0930-lengthmagnitude, def:physics:phyx-mini-0930:phyxminiproblems-problemphyxmini0930-nonnegativesireadout}
  Read a length in metres
\end{definition}
\begin{proof}
  This is the declared model definition of length In Meters.
\end{proof}
\begin{definition}[length In Centimeters]
  \label{def:physics:phyx-mini-0930:phyxminiproblems-problemphyxmini0930-lengthincentimeters}
  \lean{PhyXMiniProblems.ProblemPhyXMini0930.lengthInCentimeters}
  \uses{def:physics:phyx-mini-0930:phyxminiproblems-problemphyxmini0930-lengthmagnitude, def:physics:phyx-mini-0930:phyxminiproblems-problemphyxmini0930-lengthinmeters}
  Read a length in centimetres
\end{definition}
\begin{proof}
  This is the declared model definition of length In Centimeters.
\end{proof}
\begin{definition}[time In Seconds]
  \label{def:physics:phyx-mini-0930:phyxminiproblems-problemphyxmini0930-timeinseconds}
  \lean{PhyXMiniProblems.ProblemPhyXMini0930.timeInSeconds}
  \uses{def:physics:phyx-mini-0930:phyxminiproblems-problemphyxmini0930-timemagnitude, def:physics:phyx-mini-0930:phyxminiproblems-problemphyxmini0930-nonnegativesireadout}
  Read a physical time in seconds
\end{definition}
\begin{proof}
  This is the declared model definition of time In Seconds.
\end{proof}
\begin{definition}[area In Square Meters]
  \label{def:physics:phyx-mini-0930:phyxminiproblems-problemphyxmini0930-areainsquaremeters}
  \lean{PhyXMiniProblems.ProblemPhyXMini0930.areaInSquareMeters}
  \uses{def:physics:phyx-mini-0930:phyxminiproblems-problemphyxmini0930-areamagnitude, def:physics:phyx-mini-0930:phyxminiproblems-problemphyxmini0930-nonnegativesireadout}
  Read an area in square metres
\end{definition}
\begin{proof}
  This is the declared model definition of area In Square Meters.
\end{proof}
\begin{definition}[signed Current In Amperes]
  \label{def:physics:phyx-mini-0930:phyxminiproblems-problemphyxmini0930-signedcurrentinamperes}
  \lean{PhyXMiniProblems.ProblemPhyXMini0930.signedCurrentInAmperes}
  \uses{def:physics:phyx-mini-0930:phyxminiproblems-problemphyxmini0930-signedelectriccurrent, def:physics:phyx-mini-0930:phyxminiproblems-problemphyxmini0930-signedsireadout}
  Read a signed electric current in amperes
\end{definition}
\begin{proof}
  This is the declared model definition of signed Current In Amperes.
\end{proof}
\begin{definition}[signed Current In Microamperes]
  \label{def:physics:phyx-mini-0930:phyxminiproblems-problemphyxmini0930-signedcurrentinmicroamperes}
  \lean{PhyXMiniProblems.ProblemPhyXMini0930.signedCurrentInMicroamperes}
  \uses{def:physics:phyx-mini-0930:phyxminiproblems-problemphyxmini0930-signedelectriccurrent, def:physics:phyx-mini-0930:phyxminiproblems-problemphyxmini0930-signedcurrentinamperes}
  Read a signed electric current in microamperes
\end{definition}
\begin{proof}
  This is the declared model definition of signed Current In Microamperes.
\end{proof}
\begin{definition}[signed Current Rate In Amperes Per Second]
  \label{def:physics:phyx-mini-0930:phyxminiproblems-problemphyxmini0930-signedcurrentrateinamperespersecond}
  \lean{PhyXMiniProblems.ProblemPhyXMini0930.signedCurrentRateInAmperesPerSecond}
  \uses{def:physics:phyx-mini-0930:phyxminiproblems-problemphyxmini0930-signedelectriccurrentrate, def:physics:phyx-mini-0930:phyxminiproblems-problemphyxmini0930-signedsireadout}
  Read a signed current-change rate in amperes per second
\end{definition}
\begin{proof}
  This is the declared model definition of signed Current Rate In Amperes Per Second.
\end{proof}
\begin{definition}[magnetic Permeability In Newtons Per Ampere Squared]
  \label{def:physics:phyx-mini-0930:phyxminiproblems-problemphyxmini0930-magneticpermeabilityinnewtonsperamperesquared}
  \lean{PhyXMiniProblems.ProblemPhyXMini0930.magneticPermeabilityInNewtonsPerAmpereSquared}
  \uses{def:physics:phyx-mini-0930:phyxminiproblems-problemphyxmini0930-magneticpermeabilitymagnitude, def:physics:phyx-mini-0930:phyxminiproblems-problemphyxmini0930-nonnegativesireadout}
  Read magnetic permeability in newtons per ampere squared
\end{definition}
\begin{proof}
  This is the declared model definition of magnetic Permeability In Newtons Per Ampere Squared.
\end{proof}
\begin{definition}[signed Magnetic Flux Density In Teslas]
  \label{def:physics:phyx-mini-0930:phyxminiproblems-problemphyxmini0930-signedmagneticfluxdensityinteslas}
  \lean{PhyXMiniProblems.ProblemPhyXMini0930.signedMagneticFluxDensityInTeslas}
  \uses{def:physics:phyx-mini-0930:phyxminiproblems-problemphyxmini0930-signedmagneticfluxdensity, def:physics:phyx-mini-0930:phyxminiproblems-problemphyxmini0930-signedsireadout}
  Read a signed axial magnetic-flux-density component in teslas
\end{definition}
\begin{proof}
  This is the declared model definition of signed Magnetic Flux Density In Teslas.
\end{proof}
\begin{definition}[signed Magnetic Flux Density Rate In Teslas Per Second]
  \label{def:physics:phyx-mini-0930:phyxminiproblems-problemphyxmini0930-signedmagneticfluxdensityrateinteslaspersecond}
  \lean{PhyXMiniProblems.ProblemPhyXMini0930.signedMagneticFluxDensityRateInTeslasPerSecond}
  \uses{def:physics:phyx-mini-0930:phyxminiproblems-problemphyxmini0930-signedmagneticfluxdensityrate, def:physics:phyx-mini-0930:phyxminiproblems-problemphyxmini0930-signedsireadout}
  Read a signed flux-density change rate in teslas per second
\end{definition}
\begin{proof}
  This is the declared model definition of signed Magnetic Flux Density Rate In Teslas Per Second.
\end{proof}
\begin{definition}[signed Magnetic Flux In Webers]
  \label{def:physics:phyx-mini-0930:phyxminiproblems-problemphyxmini0930-signedmagneticfluxinwebers}
  \lean{PhyXMiniProblems.ProblemPhyXMini0930.signedMagneticFluxInWebers}
  \uses{def:physics:phyx-mini-0930:phyxminiproblems-problemphyxmini0930-signedmagneticflux, def:physics:phyx-mini-0930:phyxminiproblems-problemphyxmini0930-signedsireadout}
  Read signed magnetic flux in webers
\end{definition}
\begin{proof}
  This is the declared model definition of signed Magnetic Flux In Webers.
\end{proof}
\begin{definition}[signed Magnetic Flux Rate In Webers Per Second]
  \label{def:physics:phyx-mini-0930:phyxminiproblems-problemphyxmini0930-signedmagneticfluxrateinweberspersecond}
  \lean{PhyXMiniProblems.ProblemPhyXMini0930.signedMagneticFluxRateInWebersPerSecond}
  \uses{def:physics:phyx-mini-0930:phyxminiproblems-problemphyxmini0930-signedmagneticfluxrate, def:physics:phyx-mini-0930:phyxminiproblems-problemphyxmini0930-signedsireadout}
  Read signed magnetic-flux change rate in webers per second
\end{definition}
\begin{proof}
  This is the declared model definition of signed Magnetic Flux Rate In Webers Per Second.
\end{proof}
\begin{definition}[signed Emf In Volts]
  \label{def:physics:phyx-mini-0930:phyxminiproblems-problemphyxmini0930-signedemfinvolts}
  \lean{PhyXMiniProblems.ProblemPhyXMini0930.signedEmfInVolts}
  \uses{def:physics:phyx-mini-0930:phyxminiproblems-problemphyxmini0930-signedelectromotiveforce, def:physics:phyx-mini-0930:phyxminiproblems-problemphyxmini0930-signedsireadout}
  Read signed electromotive force in volts
\end{definition}
\begin{proof}
  This is the declared model definition of signed Emf In Volts.
\end{proof}
\begin{definition}[resistance In Ohms]
  \label{def:physics:phyx-mini-0930:phyxminiproblems-problemphyxmini0930-resistanceinohms}
  \lean{PhyXMiniProblems.ProblemPhyXMini0930.resistanceInOhms}
  \uses{def:physics:phyx-mini-0930:phyxminiproblems-problemphyxmini0930-electricalresistancemagnitude, def:physics:phyx-mini-0930:phyxminiproblems-problemphyxmini0930-nonnegativesireadout}
  Read electrical resistance in ohms
\end{definition}
\begin{proof}
  This is the declared model definition of resistance In Ohms.
\end{proof}
\begin{definition}[Circuit Sense]
  \label{def:physics:phyx-mini-0930:phyxminiproblems-problemphyxmini0930-circuitsense}
  \lean{PhyXMiniProblems.ProblemPhyXMini0930.CircuitSense}
  Sense of circulation as viewed along the line of sight
\end{definition}
\begin{proof}
  This is the declared model definition of Circuit Sense.
\end{proof}
\begin{definition}[Viewing Side]
  \label{def:physics:phyx-mini-0930:phyxminiproblems-problemphyxmini0930-viewingside}
  \lean{PhyXMiniProblems.ProblemPhyXMini0930.ViewingSide}
  The side from which the circular windings are viewed in the problem
\end{definition}
\begin{proof}
  This is the declared model definition of Viewing Side.
\end{proof}
\begin{definition}[Solenoid Axial Direction]
  \label{def:physics:phyx-mini-0930:phyxminiproblems-problemphyxmini0930-solenoidaxialdirection}
  \lean{PhyXMiniProblems.ProblemPhyXMini0930.SolenoidAxialDirection}
  Axial directions distinguished relative to the observer at the left
\end{definition}
\begin{proof}
  This is the declared model definition of Solenoid Axial Direction.
\end{proof}
\begin{definition}[Solenoid Model]
  \label{def:physics:phyx-mini-0930:phyxminiproblems-problemphyxmini0930-solenoidmodel}
  \lean{PhyXMiniProblems.ProblemPhyXMini0930.SolenoidModel}
  Idealization of the wound conductor
\end{definition}
\begin{proof}
  This is the declared model definition of Solenoid Model.
\end{proof}
\begin{definition}[Surrounding Loop Model]
  \label{def:physics:phyx-mini-0930:phyxminiproblems-problemphyxmini0930-surroundingloopmodel}
  \lean{PhyXMiniProblems.ProblemPhyXMini0930.SurroundingLoopModel}
  Idealization of the surrounding orange conducting loop
\end{definition}
\begin{proof}
  This is the declared model definition of Surrounding Loop Model.
\end{proof}
\begin{definition}[Figure Object]
  \label{def:physics:phyx-mini-0930:phyxminiproblems-problemphyxmini0930-figureobject}
  \lean{PhyXMiniProblems.ProblemPhyXMini0930.FigureObject}
  Physical and graphical objects visible in 930.png
\end{definition}
\begin{proof}
  This is the declared model definition of Figure Object.
\end{proof}
\begin{definition}[Figure Label]
  \label{def:physics:phyx-mini-0930:phyxminiproblems-problemphyxmini0930-figurelabel}
  \lean{PhyXMiniProblems.ProblemPhyXMini0930.FigureLabel}
  \uses{def:physics:phyx-mini-0930:phyxminiproblems-problemphyxmini0930-timeinseconds}
  Symbolic or dimensional labels printed in 930.png
\end{definition}
\begin{proof}
  This is the declared model definition of Figure Label.
\end{proof}
\begin{definition}[Solenoid Loop Figure]
  \label{def:physics:phyx-mini-0930:phyxminiproblems-problemphyxmini0930-solenoidloopfigure}
  \lean{PhyXMiniProblems.ProblemPhyXMini0930.SolenoidLoopFigure}
  \uses{def:physics:phyx-mini-0930:phyxminiproblems-problemphyxmini0930-circuitsense, def:physics:phyx-mini-0930:phyxminiproblems-problemphyxmini0930-viewingside, def:physics:phyx-mini-0930:phyxminiproblems-problemphyxmini0930-figureobject, def:physics:phyx-mini-0930:phyxminiproblems-problemphyxmini0930-figurelabel}
  Literal presentation data from the primary raster. Axis coordinates and printed diameter values are scalar readouts in the units named by their fields
\end{definition}
\begin{proof}
  This is the declared model definition of Solenoid Loop Figure.
\end{proof}
\begin{definition}[Solenoid Loop Setup]
  \label{def:physics:phyx-mini-0930:phyxminiproblems-problemphyxmini0930-solenoidloopsetup}
  \lean{PhyXMiniProblems.ProblemPhyXMini0930.SolenoidLoopSetup}
  \uses{def:physics:phyx-mini-0930:phyxminiproblems-problemphyxmini0930-lengthmagnitude, def:physics:phyx-mini-0930:phyxminiproblems-problemphyxmini0930-timemagnitude, def:physics:phyx-mini-0930:phyxminiproblems-problemphyxmini0930-areamagnitude, def:physics:phyx-mini-0930:phyxminiproblems-problemphyxmini0930-signedelectriccurrent, def:physics:phyx-mini-0930:phyxminiproblems-problemphyxmini0930-signedelectriccurrentrate, def:physics:phyx-mini-0930:phyxminiproblems-problemphyxmini0930-magneticpermeabilitymagnitude, def:physics:phyx-mini-0930:phyxminiproblems-problemphyxmini0930-signedmagneticfluxdensity, def:physics:phyx-mini-0930:phyxminiproblems-problemphyxmini0930-signedmagneticfluxdensityrate, def:physics:phyx-mini-0930:phyxminiproblems-problemphyxmini0930-signedmagneticflux, def:physics:phyx-mini-0930:phyxminiproblems-problemphyxmini0930-signedmagneticfluxrate, def:physics:phyx-mini-0930:phyxminiproblems-problemphyxmini0930-signedelectromotiveforce, def:physics:phyx-mini-0930:phyxminiproblems-problemphyxmini0930-electricalresistancemagnitude, def:physics:phyx-mini-0930:phyxminiproblems-problemphyxmini0930-circuitsense, def:physics:phyx-mini-0930:phyxminiproblems-problemphyxmini0930-viewingside, def:physics:phyx-mini-0930:phyxminiproblems-problemphyxmini0930-solenoidaxialdirection, def:physics:phyx-mini-0930:phyxminiproblems-problemphyxmini0930-solenoidmodel, def:physics:phyx-mini-0930:phyxminiproblems-problemphyxmini0930-surroundingloopmodel, def:physics:phyx-mini-0930:phyxminiproblems-problemphyxmini0930-solenoidloopfigure}
  Independent physical observables of the coupled solenoid-loop system. Functions ending in AtSecond take a real coordinate explicitly measured in seconds. Their values remain unit-independent physical quantities. The field component and linked flux are positive along positiveAxialDirection; emf and loop current are positive in loopPositiveCurrentSense
\end{definition}
\begin{proof}
  This is the declared model definition of Solenoid Loop Setup.
\end{proof}
\begin{definition}[Matches Stated Solenoid Loop Scenario]
  \label{def:physics:phyx-mini-0930:phyxminiproblems-problemphyxmini0930-matchesstatedsolenoidloopscenario}
  \lean{PhyXMiniProblems.ProblemPhyXMini0930.MatchesStatedSolenoidLoopScenario}
  \uses{def:physics:phyx-mini-0930:phyxminiproblems-problemphyxmini0930-lengthincentimeters, def:physics:phyx-mini-0930:phyxminiproblems-problemphyxmini0930-timeinseconds, def:physics:phyx-mini-0930:phyxminiproblems-problemphyxmini0930-resistanceinohms, def:physics:phyx-mini-0930:phyxminiproblems-problemphyxmini0930-solenoidloopsetup}
  Prose-supplied apparatus data and the requested observation time
\end{definition}
\begin{proof}
  This is the declared model definition of Matches Stated Solenoid Loop Scenario.
\end{proof}
\begin{definition}[Matches Supplied Solenoid Loop Figure]
  \label{def:physics:phyx-mini-0930:phyxminiproblems-problemphyxmini0930-matchessuppliedsolenoidloopfigure}
  \lean{PhyXMiniProblems.ProblemPhyXMini0930.MatchesSuppliedSolenoidLoopFigure}
  \uses{def:physics:phyx-mini-0930:phyxminiproblems-problemphyxmini0930-lengthincentimeters, def:physics:phyx-mini-0930:phyxminiproblems-problemphyxmini0930-timeinseconds, def:physics:phyx-mini-0930:phyxminiproblems-problemphyxmini0930-signedcurrentinamperes, def:physics:phyx-mini-0930:phyxminiproblems-problemphyxmini0930-solenoidloopsetup}
  Primary-raster geometry, labels, axes, and the complete three-piece current trace. The graph calibrates only the solenoid current; it contains no loop current or induced-emf answer
\end{definition}
\begin{proof}
  This is the declared model definition of Matches Supplied Solenoid Loop Figure.
\end{proof}
\begin{definition}[Has Physical Solenoid Loop Parameters]
  \label{def:physics:phyx-mini-0930:phyxminiproblems-problemphyxmini0930-hasphysicalsolenoidloopparameters}
  \lean{PhyXMiniProblems.ProblemPhyXMini0930.HasPhysicalSolenoidLoopParameters}
  \uses{def:physics:phyx-mini-0930:phyxminiproblems-problemphyxmini0930-lengthinmeters, def:physics:phyx-mini-0930:phyxminiproblems-problemphyxmini0930-areainsquaremeters, def:physics:phyx-mini-0930:phyxminiproblems-problemphyxmini0930-magneticpermeabilityinnewtonsperamperesquared, def:physics:phyx-mini-0930:phyxminiproblems-problemphyxmini0930-resistanceinohms, def:physics:phyx-mini-0930:phyxminiproblems-problemphyxmini0930-solenoidloopsetup}
  Positivity and geometric containment conditions for the physical setup
\end{definition}
\begin{proof}
  This is the declared model definition of Has Physical Solenoid Loop Parameters.
\end{proof}
\begin{definition}[Satisfies Circular Solenoid Geometry]
  \label{def:physics:phyx-mini-0930:phyxminiproblems-problemphyxmini0930-satisfiescircularsolenoidgeometry}
  \lean{PhyXMiniProblems.ProblemPhyXMini0930.SatisfiesCircularSolenoidGeometry}
  \uses{def:physics:phyx-mini-0930:phyxminiproblems-problemphyxmini0930-lengthinmeters, def:physics:phyx-mini-0930:phyxminiproblems-problemphyxmini0930-areainsquaremeters, def:physics:phyx-mini-0930:phyxminiproblems-problemphyxmini0930-solenoidloopsetup}
  Circular geometry fixes the area through which ideal-solenoid flux links
\end{definition}
\begin{proof}
  This is the declared model definition of Satisfies Circular Solenoid Geometry.
\end{proof}
\begin{definition}[Satisfies Solenoid Current Rate Definition]
  \label{def:physics:phyx-mini-0930:phyxminiproblems-problemphyxmini0930-satisfiessolenoidcurrentratedefinition}
  \lean{PhyXMiniProblems.ProblemPhyXMini0930.SatisfiesSolenoidCurrentRateDefinition}
  \uses{def:physics:phyx-mini-0930:phyxminiproblems-problemphyxmini0930-timeinseconds, def:physics:phyx-mini-0930:phyxminiproblems-problemphyxmini0930-signedcurrentinamperes, def:physics:phyx-mini-0930:phyxminiproblems-problemphyxmini0930-signedcurrentrateinamperespersecond, def:physics:phyx-mini-0930:phyxminiproblems-problemphyxmini0930-solenoidloopsetup}
  The rate observable is the derivative of the graphed signed current away from the two corners of the piecewise-linear trace
\end{definition}
\begin{proof}
  This is the declared model definition of Satisfies Solenoid Current Rate Definition.
\end{proof}
\begin{definition}[Satisfies Ideal Solenoid Field Law]
  \label{def:physics:phyx-mini-0930:phyxminiproblems-problemphyxmini0930-satisfiesidealsolenoidfieldlaw}
  \lean{PhyXMiniProblems.ProblemPhyXMini0930.SatisfiesIdealSolenoidFieldLaw}
  \uses{def:physics:phyx-mini-0930:phyxminiproblems-problemphyxmini0930-lengthinmeters, def:physics:phyx-mini-0930:phyxminiproblems-problemphyxmini0930-timeinseconds, def:physics:phyx-mini-0930:phyxminiproblems-problemphyxmini0930-signedcurrentinamperes, def:physics:phyx-mini-0930:phyxminiproblems-problemphyxmini0930-signedcurrentrateinamperespersecond, def:physics:phyx-mini-0930:phyxminiproblems-problemphyxmini0930-magneticpermeabilityinnewtonsperamperesquared, def:physics:phyx-mini-0930:phyxminiproblems-problemphyxmini0930-signedmagneticfluxdensityinteslas, def:physics:phyx-mini-0930:phyxminiproblems-problemphyxmini0930-signedmagneticfluxdensityrateinteslaspersecond, def:physics:phyx-mini-0930:phyxminiproblems-problemphyxmini0930-solenoidloopsetup}
  The ideal long-solenoid law B = μ₀ (N / ℓ) I and its rate form. These laws mention neither the surrounding-loop current nor any displayed answer value
\end{definition}
\begin{proof}
  This is the declared model definition of Satisfies Ideal Solenoid Field Law.
\end{proof}
\begin{definition}[Satisfies Uniform Solenoid Flux Linkage Law]
  \label{def:physics:phyx-mini-0930:phyxminiproblems-problemphyxmini0930-satisfiesuniformsolenoidfluxlinkagelaw}
  \lean{PhyXMiniProblems.ProblemPhyXMini0930.SatisfiesUniformSolenoidFluxLinkageLaw}
  \uses{def:physics:phyx-mini-0930:phyxminiproblems-problemphyxmini0930-timeinseconds, def:physics:phyx-mini-0930:phyxminiproblems-problemphyxmini0930-areainsquaremeters, def:physics:phyx-mini-0930:phyxminiproblems-problemphyxmini0930-signedmagneticfluxdensityinteslas, def:physics:phyx-mini-0930:phyxminiproblems-problemphyxmini0930-signedmagneticfluxdensityrateinteslaspersecond, def:physics:phyx-mini-0930:phyxminiproblems-problemphyxmini0930-signedmagneticfluxinwebers, def:physics:phyx-mini-0930:phyxminiproblems-problemphyxmini0930-signedmagneticfluxrateinweberspersecond, def:physics:phyx-mini-0930:phyxminiproblems-problemphyxmini0930-solenoidloopsetup}
  The larger loop links the field only over the solenoid cross-section. Thus flux and flux rate are the signed axial field component and its rate times the same physical cross-sectional area
\end{definition}
\begin{proof}
  This is the declared model definition of Satisfies Uniform Solenoid Flux Linkage Law.
\end{proof}
\begin{definition}[Satisfies Faraday Induction Law]
  \label{def:physics:phyx-mini-0930:phyxminiproblems-problemphyxmini0930-satisfiesfaradayinductionlaw}
  \lean{PhyXMiniProblems.ProblemPhyXMini0930.SatisfiesFaradayInductionLaw}
  \uses{def:physics:phyx-mini-0930:phyxminiproblems-problemphyxmini0930-timeinseconds, def:physics:phyx-mini-0930:phyxminiproblems-problemphyxmini0930-signedmagneticfluxrateinweberspersecond, def:physics:phyx-mini-0930:phyxminiproblems-problemphyxmini0930-signedemfinvolts, def:physics:phyx-mini-0930:phyxminiproblems-problemphyxmini0930-solenoidloopsetup}
  Faraday's induction law in the selected positive boundary orientation
\end{definition}
\begin{proof}
  This is the declared model definition of Satisfies Faraday Induction Law.
\end{proof}
\begin{definition}[Satisfies Resistive Loop Ohms Law]
  \label{def:physics:phyx-mini-0930:phyxminiproblems-problemphyxmini0930-satisfiesresistiveloopohmslaw}
  \lean{PhyXMiniProblems.ProblemPhyXMini0930.SatisfiesResistiveLoopOhmsLaw}
  \uses{def:physics:phyx-mini-0930:phyxminiproblems-problemphyxmini0930-timeinseconds, def:physics:phyx-mini-0930:phyxminiproblems-problemphyxmini0930-signedcurrentinamperes, def:physics:phyx-mini-0930:phyxminiproblems-problemphyxmini0930-signedemfinvolts, def:physics:phyx-mini-0930:phyxminiproblems-problemphyxmini0930-resistanceinohms, def:physics:phyx-mini-0930:phyxminiproblems-problemphyxmini0930-solenoidloopsetup}
  Signed Ohm's law for the passive surrounding loop
\end{definition}
\begin{proof}
  This is the declared model definition of Satisfies Resistive Loop Ohms Law.
\end{proof}
\begin{definition}[Answer Choice]
  \label{def:physics:phyx-mini-0930:phyxminiproblems-problemphyxmini0930-answerchoice}
  \lean{PhyXMiniProblems.ProblemPhyXMini0930.AnswerChoice}
  Labels of the four current-valued answer choices
\end{definition}
\begin{proof}
  This is the declared model definition of Answer Choice.
\end{proof}
\begin{definition}[displayed Current In Microamperes]
  \label{def:physics:phyx-mini-0930:phyxminiproblems-problemphyxmini0930-answerchoice-displayedcurrentinmicroamperes}
  \lean{PhyXMiniProblems.ProblemPhyXMini0930.AnswerChoice.displayedCurrentInMicroamperes}
  \uses{def:physics:phyx-mini-0930:phyxminiproblems-problemphyxmini0930-answerchoice}
  Current in microamperes printed beside each answer choice
\end{definition}
\begin{proof}
  This is the declared model definition of displayed Current In Microamperes.
\end{proof}
\begin{definition}[recorded Dataset Answer]
  \label{def:physics:phyx-mini-0930:phyxminiproblems-problemphyxmini0930-recordeddatasetanswer}
  \lean{PhyXMiniProblems.ProblemPhyXMini0930.recordedDatasetAnswer}
  \uses{def:physics:phyx-mini-0930:phyxminiproblems-problemphyxmini0930-answerchoice}
  Dataset-recorded answer label
\end{definition}
\begin{proof}
  This is the declared model definition of recorded Dataset Answer.
\end{proof}
\begin{definition}[observed Loop Current In Microamperes]
  \label{def:physics:phyx-mini-0930:phyxminiproblems-problemphyxmini0930-observedloopcurrentinmicroamperes}
  \lean{PhyXMiniProblems.ProblemPhyXMini0930.observedLoopCurrentInMicroamperes}
  \uses{def:physics:phyx-mini-0930:phyxminiproblems-problemphyxmini0930-timeinseconds, def:physics:phyx-mini-0930:phyxminiproblems-problemphyxmini0930-signedcurrentinmicroamperes, def:physics:phyx-mini-0930:phyxminiproblems-problemphyxmini0930-solenoidloopsetup}
  Signed loop-current readout at the independently modeled observation time
\end{definition}
\begin{proof}
  This is the declared model definition of observed Loop Current In Microamperes.
\end{proof}
\begin{definition}[Answer Matches Observed Loop Current]
  \label{def:physics:phyx-mini-0930:phyxminiproblems-problemphyxmini0930-answermatchesobservedloopcurrent}
  \lean{PhyXMiniProblems.ProblemPhyXMini0930.AnswerMatchesObservedLoopCurrent}
  \uses{def:physics:phyx-mini-0930:phyxminiproblems-problemphyxmini0930-solenoidloopsetup, def:physics:phyx-mini-0930:phyxminiproblems-problemphyxmini0930-answerchoice, def:physics:phyx-mini-0930:phyxminiproblems-problemphyxmini0930-observedloopcurrentinmicroamperes}
  A displayed choice agrees exactly with the requested current readout
\end{definition}
\begin{proof}
  This is the declared model definition of Answer Matches Observed Loop Current.
\end{proof}
\begin{definition}[Is Unique Matching Answer]
  \label{def:physics:phyx-mini-0930:phyxminiproblems-problemphyxmini0930-isuniquematchinganswer}
  \lean{PhyXMiniProblems.ProblemPhyXMini0930.IsUniqueMatchingAnswer}
  \uses{def:physics:phyx-mini-0930:phyxminiproblems-problemphyxmini0930-solenoidloopsetup, def:physics:phyx-mini-0930:phyxminiproblems-problemphyxmini0930-answerchoice, def:physics:phyx-mini-0930:phyxminiproblems-problemphyxmini0930-answermatchesobservedloopcurrent}
  A choice is the unique displayed answer matching the requested current
\end{definition}
\begin{proof}
  This is the declared model definition of Is Unique Matching Answer.
\end{proof}
% --- Archon declaration topology end ---
% --- Archon physics formalization source end ---
